\glossentry{A priori hypothesis}{先验假设}{用于指在分析数据之前，最好也在收集数据之前，就已明确规定的假设。例如，研究者在经费申请书中列出两项先验假设：第一，在45--64岁男性中，煤尘暴露是前列腺癌的危险因素；第二，在丈夫年龄为45--64岁的女性中，配偶接触煤尘是子宫内膜癌的危险因素。}
\glossentry{Absolute risk reduction}{绝对风险降低}{治疗使结局风险降低的量，定义为未治疗者的风险减去治疗者的风险；其倒数为\emph{需治疗人数}。如果治疗反而增加风险，风险降低量就是负值，此时直接改变符号，称作绝对风险增加更为容易。绝对风险增加的倒数有时称为“需伤害人数”。例如，妇女、婴幼儿特殊补充食品计划（WIC）将低出生体重风险从10\%降至6\%，绝对风险降低为4\%。（Buescher PA, Larson LC, Nelson MD, Lenihan AJ. Prenatal WIC participation can reduce low birth weight and newborn medical costs: a cost-benefit analysis of WIC participation in North Carolina. J Acad Nutr Dietetics. 1993;93:163-166.）}
\glossentry{Accessible population}{可接触人群}{研究者能够接触、能够从中选择人员并邀请其参加研究的一群人。例如，某研究的可接触人群为2021年1月1日至2024年6月30日在 Longview 医院接受治疗的乳腺癌女性。另见\emph{预定样本、目标人群}。}
\glossentry{Accuracy}{准确度}{测量与其真实值相符合的程度，受是否存在偏倚以及精密度的影响。例如，自报体重对实际体重的测量，不如经过校准的电子秤准确。}
\glossentry{Active control trial (comparative efficacy trial)}{活性对照试验（比较疗效试验）}{对照组接受某种已知或被认为会影响目标结局的干预的临床试验。例如，研究者采用活性对照设计，将治疗关节痛的新药与标准剂量的非甾体抗炎药比较。
\sourcepage{671}}
\glossentry{Adaptive clinical trial design}{适应性临床试验设计}{允许根据期中监测和结果分析，修改试验方案的设计。例如，某适应性设计允许在期中分析显示干预可能有效、但可能需要更大样本才有足够检验效能时，增加样本量并延长招募时间。}
\glossentry{Adaptive randomization}{适应性随机化}{旨在保证研究组间指定协变量平衡的随机化技术。每名合格参与者的随机分配，会根据已经入组者的特征作相应调整。例如，在预防哮喘住院的新治疗研究中，可用此法使既往曾需口服皮质激素者，更有可能被分到其入组当时平均住院风险较低的组。}
\glossentry{Adjustment}{调整}{各种统计技术的总称，用于考虑一个或多个变量对另外两个变量之间关联的影响。例如，调整母亲吸烟后，母亲吸大麻与出生体重之间的关联幅度减小。另见\emph{条件化}。}
\glossentry{Adverse events}{不良事件}{研究期间发生、对健康或福祉产生不利影响的事件，通常用于临床试验。干预组和对照组都记录不良事件；事件不必与干预有关。例如，某研究最常见的不良事件是上呼吸道症状。}
\glossentry{Alpha}{α}{设计研究时，预先规定犯I类错误的最大概率，即零假设为真时却拒绝它的概率。例如，研究者选择$\alpha=0.05$，就是规定：若目标人群中非白人种族与结肠癌风险本无关联，研究仅因偶然而发现统计学显著关联的概率上限为5\%。也称\emph{统计学显著性水平}。}
\glossentry{Alternative hypothesis}{备择假设}{估计样本量时使用的命题，认为人群中预测变量与结局变量之间存在关联。例如，研究的备择假设是：吸烟青少年与不吸烟者辍学的可能性不同。另见\emph{零假设}。
\sourcepage{672}}
\glossentry{Analytic study}{分析性研究}{寻找两个或多个变量之间关联的研究。例如，研究者开展分析性研究，考察医学生的身高是否与血压相关。另见\emph{描述性研究}。}
\glossentry{Ancillary study}{附属研究}{在另一项原先未设计该测量的研究中，增加一个或多个预测因素或结局的研究。例如，某炎症性肠病新药随机试验保存了参与者在研究开始时采集的血清。附属研究可使用这些血清，检验某种生物标志物能否预测对药物的反应。}
\glossentry{Association}{关联}{两个变量之间可量化的关系。例如，研究发现60--69岁人群中男性与认知障碍风险相关，风险比为1.3。}
\glossentry{As-treated analysis}{按实际治疗分析}{分析临床试验资料时，按参与者实际接受的治疗（干预或对照），而非原随机分配的治疗来比较结局。例如，一项儿童复发中耳炎试验比较手术（鼓膜切开、放置引流管）与药物治疗（抗生素）。按实际治疗分析时，随机分到药物组但接受手术者归入手术组；随机分到手术组但未接受手术者归入药物组。另见\emph{符合方案分析、意向治疗分析}。}
\glossentry{Bayesian analysis}{贝叶斯分析}{得名于 Thomas Bayes 牧师的数据分析方法。根据研究结果进行推断时，它使用关于研究假设可能性的累积知识。例如，纳入类似操作既往试验结果的贝叶斯分析认为，某种新手术有95\%的可能性，将早期胰腺癌患者5年死亡率降低至少10\%。另见\emph{频率学派分析}。}
\glossentry{Bayesian trial}{贝叶斯试验}{先对治疗效果设定先验分布，再用研究结果修正，形成后验分布的试验。它可以没有固定样本量，而持续到治疗有效、无效或两者皆非的可能性足够明确。例如，既往治疗试验显示心血管结局风险降低30\%，但存在不确定性，如仍有10\%概率治疗无效。后续同治疗的贝叶斯试验纳入足够人数后，若发现风险降低25\%，可能得到后验判断：治疗有99\%概率使风险至少降低20\%，有90\%概率使风险至少降低25\%。
\sourcepage{673}}
\glossentry{Before-after study (or trial) design}{前后研究（试验）设计}{比较实施干预之前与之后结局的研究，有时称干预前后（pre-post）设计。例如，研究者比较 ICU 在实施减少不恰当导尿的措施前后，导管相关尿路感染的发生率。}
\glossentry{Beneficence}{有利原则}{研究伦理的基本原则，要求研究所得科学知识的价值超过参与者承受的不便和风险，并尽量降低风险。例如，要求参与者接受风险很高的干预，如在疫苗挑战试验中暴露于 COVID-19 病毒，可能违反有利原则。}
\glossentry{Beta}{β}{设计时预设犯II类错误的最大概率，即零假设为假却未能拒绝它的概率。只有结合效应量，这一指标才有意义。例如，结肠癌少见（中年人每年每10万人约20例）；若规定$\beta=0.20$、双侧$\alpha=0.05$，要证明每日阿司匹林相较安慰剂使10年结肠癌风险减半，每组约需25,000人。换言之，即使阿司匹林效果恰好如此，这样的研究仍有20\%概率不能在$\alpha=0.05$水平拒绝无差异的零假设。另见\emph{检验效能}。}
\glossentry{Between-groups design}{组间设计}{比较两个或多个不同参与者群体的特征或结局的设计。例如，研究者比较两类 ICU 患者的院内死亡率：一类全天有重症专科医生在现场，另一类由中央枢纽的重症医生通过电子系统监测。另见\emph{组内设计}。}
\glossentry{Bias}{偏倚}{研究设计、执行或分析中的不足，导致测量或关联估计出现系统误差。例如，由于回忆有毒化学品暴露的方式存在偏倚，白血病患者比对照更可能报告过去使用过杀虫剂。
\sourcepage{674}}
\glossentry{Biomarker}{生物标志物}{生物过程的一项客观、可量化特征。例如，血清 CA-125 水平常用作卵巢癌女性随访中的生物标志物。另见\emph{中间标志物}。}
\glossentry{Blinding}{盲法}{保证参与者、治疗医生和／或研究者不知道参与者被分入哪组（如干预或对照）的过程，通常用于随机试验。也称 masking，尤其见于眼科研究。例如，使用外观相同的安慰剂片，并将分组名单存放在异地，使参与者和研究者（包括研究助理）都不知道谁接受活性药。}
\glossentry{Blocked randomization}{区组随机化}{按预定大小（如4人或6人）的区组分配干预，以确保干预与对照组人数接近。多中心希望各场所两组总人数相近时常用。例如，每个诊所以6人为区组随机分到治疗或对照，保证任何时候两组人数相差不超过3。另见\emph{分层区组随机化}。}
\glossentry{Bonferroni correction}{Bonferroni校正}{将研究的总体$\alpha$除以所检验的假设数，减少I类错误的方法。例如，研究者检验四项不同假设，因而把每项$\alpha$从0.05降至0.0125。}
\glossentry{Bootstrapping}{自助法}{反复从一个数据集中有放回地随机抽样，估计目标参数的方法。每次样本得到的参数估计略有不同，因此可构造置信区间及其他估计变异指标。例如，用后退逐步 logistic 回归为临床预测模型选变量的研究者，可在多次自助样本中重复该过程，估计所得规则灵敏度和特异度的95\%置信区间。}
\glossentry{Calibration}{校准}{（1）确保仪器提供一致读数的过程，通常测量已知标准，再据此调整仪器。例如，每月用3 kg钢块校准新生儿体重秤。（2）事件预测概率与实际观察概率相符合的程度，常以横轴预测、纵轴实际概率作图。例如，美国心脏病学会合并队列方程常用于估计10年心血管风险，指导他汀决策；但对较近年代的队列，校准可能不佳，高估未来心血管事件风险。
\sourcepage{675}}
\glossentry{Carryover effect}{残留效应}{停止干预后仍残留的影响（如果存在）。例如，由于残留效应，双膦酸盐疗程结束后，骨密度可能多年仍未回到治疗开始前的水平。}
\glossentry{Case}{病例}{已经发生或之后发生目标结局的参与者。例如，将2018--2022年在圣迭戈县诊断恶性黑色素瘤者定义为病例。另见\emph{对照}。}
\glossentry{Case series}{病例系列}{研究一系列被诊断为特定疾病者，有时他们还接受了特定治疗，通常用于描述或产生假设。例如，医师助理实施白内障手术患者的病例系列，报告所有患者结局满意。}
\glossentry{Case-cohort study}{病例队列研究}{在较大队列随访中，选取发生疾病或其他结局者为病例，与整个队列的随机样本比较。例如，在2000名早期前列腺癌男性队列中，比较随访死于前列腺癌者与全队列随机样本在研究开始时采样测得的雄激素、维生素D水平。}
\glossentry{Case-control study}{病例对照研究}{将已有或发生疾病或其他结局的病例，与没有该结局的对照比较的设计。例如，比较急诊所见憩室炎病例与其他胃肠诊断对照，每周平均坚果和种子摄入量。}
\glossentry{Case-crossover study}{病例交叉研究}{一种病例对照设计，每个病例自身作对照，将结局发生前某时期的时变暴露，与一个或多个对照时期比较。设计易受回忆偏倚，故暴露能客观判定时最有用。例如，研究急诊肺栓塞患者是否在前48小时比前一周相同时段更可能乘过飞机。
\sourcepage{676}}
\glossentry{Categorical variable}{分类变量}{只能取若干可能值之一的变量。例如，把自报教育程度转为四类：高中以下、高中或读过部分大学、仅大学学位、接受过任何研究生教育。另见\emph{连续变量、二分类变量、无序分类变量、有序分类变量}。}
\glossentry{Causal effect}{因果效应}{对个体层面的二分类暴露与结局，反事实地比较：此人有目标暴露时是否发生结局，与若无暴露本会怎样；最好通过随机盲法试验估计。人群层面则为全部个体效应的平均。例如，随机一半服阿司匹林、一半服匹配安慰剂，随访20年并定期检查皮肤，估计对黑色素瘤风险的平均因果效应。分类和连续变量的效应也可类似定义。}
\glossentry{Cause-effect}{因果关系}{预测变量导致结局或提高其发生可能性的概念。许多观察研究旨在证明因果关系，但除非原因（如治疗）被随机分配，否则很难做到。例如，研究者用多变量分析控制潜在混杂，考察饮酒（因）与胰腺癌（果）是否有因果关系。另见\emph{因果效应、混杂、果因关系}。}
\glossentry{Censoring}{删失}{纵向研究中，因失访、不再处于风险中，或研究结束前未发生目标事件，结局信息不完整。被删失者事件或生存时间未知，只知道大于其随访时间。例如，子宫癌危险因素研究在参与者接受子宫切除时删失。}
\glossentry{Certificate of confidentiality}{保密证书}{联邦机构签发的文件，禁止未经参与者同意，向任何与研究无关的人披露可识别个人的资料，包括因传票或法院命令而披露。例如，保密证书保护苯丙胺使用研究的研究者，使其不必向执法机构提供数据。
\sourcepage{677}}
\glossentry{Chi-squared test}{卡方检验}{比较两个或多个比例，判断是否显著不同的统计检验。例如，研究比较每周至少运动两次者与运动较少者的痴呆风险，用卡方检验判断风险是否相似。}
\glossentry{Classification and Regression Tree (CART)}{分类与回归树}{见\emph{递归分割}。}
\glossentry{Clinic-based control}{诊所来源对照}{病例对照研究中，从病例来源的同一诊所或诊疗机构选出的对照。例如，研究每周至少在硬路面跑2英里是否与膝骨关节炎影像表现相关，采用同诊所对照。}
\glossentry{Clinical data managers}{临床数据管理员}{创建录入表、管理监测收集过程、整理并提取分析数据的专门人员，通常服务于为药物器械监管审批作准备的试验。例如，临床数据管理员可能因举报所发现的研究不当行为，而成为百万富翁。}
\glossentry{Clinical prediction model}{临床预测模型}{结合多个预测因素，包括各种症状体征有无及医学检查结果，估计疾病或结局概率的算法。例如，依据既往骨折、跌倒特征（如果发生过）和前臂体检，开发绝经后女性腕骨折诊断模型。}
\glossentry{Clinical trial}{临床试验}{参与者接受至少两种不同干预之一的设计。通常且最好随机分配，因此称随机临床试验。例如，研究异常心脏瓣膜患者接受牙科操作时，预防性青霉素能否降低细菌性心内膜炎风险。另见\emph{随机盲法试验}。}
\glossentry{Closed-ended question}{封闭式问题}{要求受访者从两个或多个预定选项中选择的问题。例如，一名自然神论者不知道宗教偏好题如何回答，因为清单没有该选项。另见\emph{开放式问题}。}
\sourcepage{678}
\glossentry{Close-out}{结束阶段}{参与者完成最后访视，或资料收集终止的时期。例如，结束程序规定最后访视后，应向所有人寄送研究期间完成的全部化验结果。}
\glossentry{Cluster randomization}{整群随机化}{把参与者群组（群集）而非每个个人随机分配到不同治疗的技术。例如，考察减少噪声对心脏术后恢复的作用，将40所医院 ICU 随机分到“术后安静”干预或“常规照护”对照。}
\glossentry{Cluster sampling}{整群抽样}{以群组而非个人选择参与者，常为方便大人群抽样而采用。例如，药物滥用患病率研究先在某电话区号内选10个三位前缀（如285、336），再在每个前缀群内随机拨号招募30名志愿者，共300人。}
\glossentry{Codebook}{编码手册}{定性研究中形成的主题及其他编码清单，包含每个编码的定义和概念关系，常以层级或树状结构呈现。例如，医患沟通编码包括 Com\_Active、Com\_Disengaged、Com\_GoWithTheFlow，每项均附定义和示例，说明怎样使用。}
\glossentry{Coding}{编码}{见\emph{主题编码}。}
\glossentry{Coefficient of variation (CV)}{变异系数}{衡量测量精密度的指标，通过同一样本一系列重复测量的均值除以这些测量的标准差得到。\footnote{译注：此处底本把分子分母反写；本书第4章原PDF84页给出的定义是标准差除以均值。}有时分别在测量范围中部和极端水平求 CV。例如，某围绝经期女性雌二醇很低，实验室 CV 为10\%，年轻女性则仅2\%。}
\glossentry{Cohort study}{队列研究}{前瞻性队列纳入一群参与者，进行初始测量，再向前随访结局；回顾性队列识别一群已经测量且已经完成随访者。有些同时具有两者特征。例如，考察美国陆军士兵入伍时的情商测试结果与发生创伤后应激障碍（PTSD）的可能性是否相关，既包括研究开始前服役期间诊断的 PTSD，也包括后续随访期间的诊断。
\sourcepage{679}}
\glossentry{Cointervention}{共同干预}{临床试验中，在随机化之后发生、影响结局可能性、但并非所研究干预的另一项干预。若两组发生率不同，会偏倚结局，使人难以把因果作用归于目标干预。例如，比较颈动脉支架与手术对后续卒中的作用时，支架组还更可能长期接受抗凝，结果因而难以解释。}
\glossentry{Collider}{碰撞点}{见\emph{对共同结果条件化}。}
\glossentry{Common data models}{共同数据模型}{规定一个或一组数据表的结构，常采用标准变量名、连续变量单位和分类变量值集，使不同来源能够合并分析。例如，OMRON 与 iHealth 家庭监测系统的收缩压可能分别用 sbp 和 sysbp，日期时间为 date＝“12/3/21 03:05:44”和 dt＝03dec2021。统一为共同血压模型，如 SBP、date＝“12/03/2021”，便于跨来源计算，但可能丢失一天中具体时刻等信息。}
\glossentry{Community advisory board}{社区咨询委员会}{在研究早期召集不同领域社区代表，持续对项目许多方面提供意见反馈。例如，一项孟加拉裔社区心血管健康社会决定因素研究的委员会，将建议收集哪些测量、怎样促进参与和招募、如何向社区传播文化相关结果，以及当前研究可能衍生的未来问题。}
\glossentry{Community engagement studio}{社区参与工作坊}{由社区利益相关者组成，针对项目具体方面及时反馈的小组。例如，研究改善美国越南裔人群乳腺癌筛查的干预时，可借助工作坊中越南裔社区的不同成员，包括乳腺癌患者及其照护者、医务人员和其他当地领导者，对拟议研究的设计和实施提出意见。
\sourcepage{680}}
\glossentry{Community outreach}{社区外展}{常见的有限强度参与策略，向社区说明拟议研究的需要和理由，常帮助招募与留存。例如，研究者向老年移民社区作新痴呆预防研究的教育讲座，提高对问题的认识。}
\glossentry{Community-based participatory research}{社区本位参与式研究}{社区参与式研究的一种基本方法，社区利益相关者参与研究所有环节，学术与社区伙伴共享权力和治理。例如，旧金山亚裔美国人癌症意识、研究与培训项目，是学术研究者与社区伙伴长期合作促进癌症认识、预防的项目。}
\glossentry{Community-engaged research}{社区参与式研究}{按照社区需求设计项目，由学术研究者和社区利益相关者以不同程度合作的研究取向。例如，如何改善筛查率低的社区乳腺癌筛查，可从当地知识和意见获益。}
\glossentry{Comparative efficacy trial}{比较疗效试验}{见\emph{活性对照试验}。}
\glossentry{Comparative study}{比较性研究}{定性研究中，既寻找场所或参与者相似之处，也识别重要差异的研究。例如，比较一个工作场所体重管理项目为何在一处纳入近半合格人员，而另一处不足10\%。另见\emph{描述性研究、探索性研究}。}
\glossentry{Complex hypothesis}{复杂假设}{含不止一个预测或结局变量的研究假设。由于难以统计检验，应避免使用。例如，把“新病例管理项目会影响住院时间和再入院可能性”的复杂假设，拆成“影响住院时间”及“影响再入院可能性”两个简单假设。
\sourcepage{681}}
\glossentry{Composite outcome}{复合结局}{由多个相关事件或测量构成的结局，通常计研究中任一事件首次发生。例如，心血管复合结局可包括心血管死亡，或因心梗、卒中、冠脉再血管化、心衰而住院。}
\glossentry{Conceptual adequacy}{概念适当性}{测量工具适当且全面反映目标概念或特征的程度。例如，研究者通过判断一小组受访者是否按预期理解其含义，测试新工具的概念适当性。}
\glossentry{Conceptual equivalence}{概念等价性}{某特征或概念的测量，在不同人群中含义相似的程度。例如，抑郁问卷预试发现，墨西哥裔美国人因对蓝色有积极联想，而不同地理解“feeling blue”，提示概念等价性可能有问题。}
\glossentry{Concordance}{一致率}{两个或多个观察者，对某现象是否发生完全一致的程度。例如，放射科医生A、B对肺叶浸润存在与否的一致率96\%，对心脏扩大则仅76\%。另见\emph{kappa}。}
\glossentry{Conditioning}{条件化}{将“被条件化”变量固定在某水平，考察另外两个或多个变量之间关联的过程。也可指研究选择标准的制定方式，见下一词条。限定、匹配、分层和多变量调整，是常见方式。例如，对过去6个月性伴人数条件化后，研究者未发现可卡因使用与梅毒风险有关联。}
\glossentry{Conditioning on a shared effect}{对共同结果条件化}{对两个变量共同导致的结果进行条件化，从而在它们之间引入关联（常为负关联）的一种偏倚；该共同结果也称碰撞点。偏倚可由入组人员选择或分析方式产生。例如，在平均每天屏幕时间至少6小时的儿童中，对总屏幕时间条件化后，电视时间越多，电子游戏时间越少；但这并不表示多看电视会减少玩游戏！
\sourcepage{682}}
\glossentry{Confidence interval}{置信区间}{这一术语常被误解。置信区间（CI）最好视为精度指标：越窄，估计越精确。例如，同为风险比3.2，95\%区间2.9--3.5比1.5--6.8精确得多。它与统计学显著性密切相关：95\%区间近似包含在$1-95\%=0.05$水平，与观察值不显著不同的各值。例如，风险比3.2、95\%区间0.9--11，在$\alpha=0.05$不显著，因为包含无效值1.0。

人们常错误地把 CI 直接解释为后验概率陈述，如真值有95\%概率位于95\%区间内。这不正确，因为后验概率还依赖研究结果之外的信息。原文进一步称，人们常把 CI 误认为条形（各值同样可能），而其“实际上呈钟形”，靠近两端的值较不可能。\footnote{译注：置信区间是两个端点限定的数值集合，本身没有钟形或条形概率密度。钟形等描述属于具体采用的分布，不能据此给区间内固定未知参数赋予后验概率。}另见\emph{α、P值}。}
\glossentry{Conflicts of interest}{利益冲突}{研究者首要利益应是有效回答重要科学问题并保护参与者安全。其他利益与这些目标冲突，可能使研究偏倚、损害客观性或削弱公众信任时，即产生利益冲突。例如，为开发新焦虑测量设备的公司担任科学顾问者，若受邀审评另一设备的论文，就存在冲突。}
\glossentry{Confounder}{混杂因素}{见\emph{混杂}。}
\glossentry{Confounding}{混杂}{预测与结局的测量关联，因为另一个同时导致二者（或位于通向二者的因果路径上）的变量，而不同于它们的因果效应的流行病学现象；该变量称混杂因素。例如，压疮与院内死亡的表面关联，部分被基础健康混杂，因为压疮患者基础健康较差。另见\emph{有向无环图、效应修饰}。}
\glossentry{Confounding by indication}{适应证混杂}{治疗适应证之一就是混杂因素的特殊混杂，通常发生在治疗与结局的观察性研究。例如，双相障碍新药与自杀风险增加的报告，可能由于基础病更重的患者被选择性地给予新药。
\sourcepage{683}}
\glossentry{Confounding variable}{混杂变量}{见\emph{混杂}。}
\glossentry{Consecutive sample}{连续样本}{逐个选择参与者，直到达到样本量的样本。通常指预定样本；病历审阅可能不需知情同意，因此也可指实际样本。例如，从2020年1月15日起，审阅风湿门诊最先100名类风湿关节炎患者的病历。}
\glossentry{Construct validity}{构念效度}{测量符合所测特征（构念）理论定义的程度。例如，朋友称“爱玩、外向”的人与被称“害羞、不太会参加聚会”的人，社会焦虑测量值差别很大，因此认为具有构念效度。另见\emph{内容效度、效标关联效度}。}
\glossentry{Contamination}{沾染}{干预的一部分或大部分效应也影响到对照组参与者的非期望过程。例如，教儿童倒数以改善总体算术能力的研究受到沾染，因为干预组忍不住把技巧教给对照组朋友。这种情形采用整群随机化更好。}
\glossentry{Content validity}{内容效度}{测量代表所研究现象多个方面的程度。例如，失眠工具测量总睡眠时间、夜间醒来、清晨早醒、起床精力及白天嗜睡，因此认为具有内容效度。另见\emph{构念效度、效标关联效度}。}
\glossentry{Continuous variable}{连续变量}{理论上可有无限多个可能值的测量。实际上，常指有“很多”可能值的数值测量，有人认为至少10个，有人认为至少20个。例如，用水银血压计以 mm Hg 测得的收缩压作为连续变量。另见\emph{分类变量、二分类变量、离散变量}。
\sourcepage{684}}
\glossentry{Control}{对照；控制}{（1）没有目标结局，因而属于与病例比较的组的参与者。例如，消化性溃疡危险因素研究，从研究同期因非胃肠诊断住院者中选对照。另见\emph{病例}。（2）试验中未接受目标干预者得到的无活性治疗（如安慰剂）或比较治疗（如常规照护）；也指接受这种替代治疗的人。例如，对照服外观与活性药完全相同的安慰剂。另见\emph{干预}。（3）作动词指试图防止非期望情形。例如，在考察参加课外活动对青少年自杀风险因果作用的观察性研究中，用 logistic 回归控制年龄、种族族裔、性别及非法药物使用等可能混杂。另见\emph{调整、条件化}。}
\glossentry{Controlled before-after study}{有对照的前后研究}{同时测量两组参与者干预前后结局的设计。例如，比较制定与未制定儿童安全座椅法的州，在立法前后车祸儿童受伤率的变化。}
\glossentry{Convenience sample}{便利样本}{因比较容易接触而选择的研究参与者。例如，研究者用本诊所患者作脑膜瘤危险因素病例对照研究的对照。}
\glossentry{Coordinating center}{协调中心}{具备组织、沟通、数据管理及统计专长的团队，确保多中心各场所达到招募目标、采用相同程序、产生可比较并可合并分析的数据。例如，协调中心从四个场所收集数据，准备分析。}
\glossentry{Correlation coefficient}{相关系数}{表示两个连续测量在线性上相关程度的统计术语，即一个改变伴随另一个成比例改变，常记$r$。例如，中年女性样本身高与体重相关，$r=0.7$。
\sourcepage{685}}
\glossentry{Count variable}{计数变量}{具有可数的数目、可表达为正整数的变量。例如，一个家庭居住人数就是计数变量。}
\glossentry{Counterfactual}{反事实}{与实际发生相反。例如，Violet 长期接触染发剂，她的反事实就是没有这种暴露。在因果推理中，通过比较实际暴露下的结局与相反（反事实）暴露下本会怎样，估计因果效应。例如，估计贫困对城市暴力犯罪的影响，可推算假设全城没有贫困人口时的犯罪率，再与假设全城人人贫困时比较。}
\glossentry{Covariate}{协变量}{除试验治疗、分析性研究主要预测变量，或任何研究结局以外的参与者特征，也称 covariable。例如，在过度接触手机屏幕是否预测自报失眠的横断面研究中，协变量包括年龄、咖啡因、酒精及镇静药使用。}
\glossentry{Cox model}{Cox模型}{也称 Cox 比例风险率模型。估计一个或多个预测变量各自对结局发生速率（风险率）的作用，并考虑参与者随访长度不同的多变量技术。例如，调整年龄、血压、糖尿病及随访时间后，男性卒中可能性约女性两倍，黑人约白人三倍。另见\emph{logistic回归模型}。}
\glossentry{Credible interval}{可信区间}{依据研究结果和给定先验，概括治疗效应后验分布的方法，常但不一定采用最高后验密度区间。例如，减重药95\%可信区间为$-1.5$至$-4.6$ kg，则药物效应位于两数之间的概率为0.95。}
\glossentry{Criterion-related validity}{效标关联效度}{测量与测同一现象的其他方式相关的程度。例如，青少年抑郁工具与 Beck 抑郁量表得分高度相关，因此认为有效标关联效度。另见\emph{构念效度、内容效度}。
\sourcepage{686}}
\glossentry{Crossover}{交叉转组}{（1）通常指试验参与者最初在一组（如常规照护），过程中转另一组（如活性治疗），活性治疗为操作时最常见。例如，原分观察等待的前列腺癌患者15\%在试验中转而接受放疗或手术。（2）更泛指任何暴露水平随时间改变。例如，急诊滑板伤病例交叉研究询问受伤当天和前一天是否看滑板视频，判断两天暴露是否发生交叉变化。}
\glossentry{Crossover study}{交叉研究}{某治疗或对照组所有参与者换至另一组的设计，通常在中点交换，有时两期之间设洗脱期。人人都能接受活性治疗，但仅适合治疗后会回到原水平的状况。例如，注意缺陷多动障碍患者参加新药与安慰剂的症状治疗交叉研究。}
\glossentry{Cross-sectional study}{横断面研究}{在一个时点或有限时期选择参与者并测量的设计，通常估计暴露或疾病患病率。例如，在加州伯克利1200名大学生中估计近视患病率。}
\glossentry{Cumulative incidence}{累积发病率}{见\emph{发病率}。}
\glossentry{DAGs}{DAG}{见\emph{有向无环图}。}
\glossentry{Data dictionary}{数据字典}{变量名及相应数据类型、格式、标签、允许值和描述的表。例如，研究者忘记 education 字段中的4表示“读过部分大学”，因此查阅字典。}
\glossentry{Data entry system}{数据录入系统}{录入研究数据的屏幕表单，以及从电子健康记录（EHR）、测量设备等来源导入资料的程序。例如，可在平板电脑完成的网络问卷，可能是录入系统的一部分。
\sourcepage{687}}
\glossentry{Data management software}{数据管理软件}{存储、管理和验证研究数据的系统。例如，研究采用所在机构免费提供、用于数据管理的 Research Electronic Data Capture（REDCap）。}
\glossentry{Data table}{数据表}{行为记录、列为字段或属性的数据表。例如，参与者表每行是一个不同参与者，各列是出生日期等不变属性。}
\glossentry{Data use agreement}{数据使用协议}{规定机构间临床数据分享的合同安排，包括分享哪些信息及协议目的。多数机构提供模板，并规定谁有权签署。例如，协议规定数据集除年龄、性别、种族族裔外不含受保护健康信息，并用于“认知障碍成人养老院安置相关因素”合作研究。}
\glossentry{Database query}{数据库查询}{用 SQL（读作“sequel”）从一表或多表选择特定记录并显示字段的指令。简单形式为 SELECT（所需字段）FROM（含这些字段的表）WHERE（限定目标记录的条件）。例如，\texttt{SELECT LastName FROM Participants WHERE FirstName = "Freddie"}只给出名叫 Freddie 的全部参与者的姓，不含其他人。}
\glossentry{Deduction}{演绎}{依据预先规定的理论或假设，自上而下考察实例的分析方法。例如，研究诊所领导如何实施新医疗服务政策，仔细考察变革管理与员工培训技术。另见\emph{归纳}。}
\glossentry{Deidentified data}{去标识化数据}{去掉可能识别参与者信息的数据。在 HIPAA 下有特定法律含义，指移除18类标识。去标识医疗资料用于研究的法律障碍少于可识别资料。例如，从当地学校系统取得数据，研究幼儿园精细运动测试表现是否影响三年级阅读理解。
\sourcepage{688}}
\glossentry{Dependent variable}{因变量}{见\emph{结局变量}。}
\glossentry{Descriptive study}{描述性研究}{（1）定量研究中，不寻找关联、不检验假设、不作比较。例如，学龄前儿童肥胖患病率研究。另见\emph{分析性研究}。（2）定性研究中，已有知识足以聚焦一个领域，但还不足以拟完整方案，例如把工具或方案调整到新目标人群。某工作场所体重管理项目在别处有效，却在此招募不足预期人数10\%，考察其原因属于例子。另见\emph{比较性研究、探索性研究}。}
\glossentry{Diagnostic test study}{诊断试验研究}{考察某医学检查结果是否影响患者特定诊断的可能性，以及影响多大的研究。例如，判断血清碳酸氢盐是否有助于在发热者中诊断脓毒症。}
\glossentry{Dichotomous variable}{二分类变量}{只能取两个值之一，如是／否、生／死的变量。例如，检查者把收缩压二分为高血压（$\ge140$ mm Hg）或非高血压。另见\emph{分类变量、连续变量}。}
\glossentry{Difference-in-differences}{双重差分}{比较干预与对照组结局变量变化之差，变化通常定义为末值减初值。例如，用此法估计病房紫外消毒相较标准清洁，对院内感染率变化的影响。}
\glossentry{Differential misclassification bias}{差异性错分偏倚}{测量随参与者状态、通常病例或对照身份而系统变化的总称，最常见于回忆暴露。例如，成人乳糜泻病例比同一家庭长大的兄弟姐妹更可能回忆童年接触含小麦食品，因此怀疑回忆偏倚，即一种差异性错分偏倚。另见\emph{非差异性错分偏倚}。
\sourcepage{689}}
\glossentry{Differential verification bias}{差异性验证偏倚}{诊断检验研究中，至少部分依据待评价检验的结果，对不同参与者采用不一定给出相同答案的不同金标准而产生的偏倚。例如，在男性前列腺癌的前列腺特异抗原（PSA）筛查研究中，PSA 水平高者接受前列腺活检，正常者则接受临床随访；这使人担心，差异性验证偏倚虚增了 PSA 筛查惰性前列腺癌的灵敏度，并压低了其特异度。}
\glossentry{Directed acyclic graph (DAG)}{有向无环图}{以示意图表示研究者认为研究变量彼此之间具有怎样的因果关系。例如，连接冒险行为、注射用药、人乳头瘤病毒（HPV）感染和宫颈癌的 DAG 可以画成：
\begin{center}\begin{tikzpicture}[>=Stealth,font=\small]
\node (risk) at (0,1.3){冒险行为};\node (drug) at (0,0){注射用药};
\node (hpv) at (4,.8){HPV感染};\node (cancer) at (6,0){宫颈癌};
\draw[->] (risk)--(drug);\draw[->] (risk)--(hpv);\draw[->] (drug)--(cancer);\draw[->] (hpv)--(cancer);
\end{tikzpicture}\end{center}
在图中，HPV感染导致宫颈癌；由于它还与注射用药共享一个原因（冒险行为），所以可能混杂注射用药与宫颈癌的表面关联。}
\glossentry{Discrete variable}{离散变量}{只取整数（可正可负）的变量，许多连续变量实际按离散变量处理。例如，过去一年体重变化记到最近的一磅。另见\emph{计数变量、连续变量}。}
\glossentry{Discrimination}{区分度}{检验或决策规则把受检人群区分为目标状况较可能和较不可能者的能力。如果有病者结果总比无病者更异常，区分就是完美的，可用受试者工作特征（ROC）曲线下面积衡量。例如，一位颅相学家声称摸法国贵族头部隆起能估计梅毒可能性，但经严格评价不优于偶然，甚至不如宫廷弄臣。}
\sourcepage{690}
\glossentry{Dose-response}{剂量反应}{暴露（剂量）越大，结局（反应）的幅度或可能性越大的现象；若暴露有保护作用，则越大结局越少。例如，日晒与黑素细胞痣数量有剂量反应。}
\glossentry{Double gold standard bias}{双重金标准偏倚}{见\emph{差异性验证偏倚}。}
\glossentry{Double-cohort study}{双队列研究}{按暴露水平不同，将参与者纳入两个不同队列，常见于同一行业不同职业。例如，比较快餐行业厨师（大部分班次站立）与收银员（坐着）的足底筋膜炎风险。}
\glossentry{Dropout}{退出者}{无法判定结局的参与者，常因拒绝随访，有时也包括研究期间死亡者。例如，失眠研究17人退出：8人拒绝继续、6人死亡、3人患痴呆。另见\emph{删失}。}
\glossentry{Ecologic fallacy}{生态学谬误}{比较人群时有关联、比较个人时却没有的情况。例如，一个国家有洗衣机的家庭比例与预期寿命相关，但个人拥有洗衣机并不会提高寿命，因此属于生态学谬误。}
\glossentry{Effect modification}{效应修饰}{第三变量（常称效应修饰因素）影响预测与结局关联强度的情况。例如，贫困对卒中风险的作用随种族不同，在黑人中的关联强于白人。它取决于所用风险模型。另见\emph{混杂、交互作用}。}
\glossentry{Effect size}{效应量}{规划样本量时，研究者希望在比较组之间检出的差异或关联大小；更泛指研究结束后该差异或关联的实际大小。例如，按两组平均血糖相差20 mg/dL估计样本量。}
\glossentry{Effect-cause}{果因关系}{结局导致预测变量，而非相反的情况。例如，病例对照研究发现吸入支气管扩张剂与间质性肺病风险增加相关，最可能是间质性肺病患者更常被错误地用吸入剂治疗。另见\emph{因果关系}。
\sourcepage{691}}
\glossentry{Effectiveness}{实际效果}{虽无标准定义，本书指干预在实际实践中的作用，区别于随机试验。例如，多个城市试验显示组织型纤溶酶原激活剂（tPA）减少卒中病残和死亡后，研究者到25所农村急诊研究其实践效果。另见\emph{疗效}。}
\glossentry{Efficacy}{疗效}{虽无标准定义，本书指干预在临床试验中的作用，区别于实际实践。例如，一项试验报告 tPA 对急性卒中患者病残与死亡的降低疗效为25\%。另见\emph{实际效果}。}
\glossentry{Entry criteria}{入组标准}{见\emph{选择标准}。}
\glossentry{Epidemiologist}{流行病学家}{一位被年龄和“性”折腾垮的临床研究者。\footnote{译注：原文“broken down by age and sex”也可读作“按年龄和性别分层”，此处故意双关。}例如，本书一位或多位作者，但我们不说是哪位！}
\glossentry{Epidemiology}{流行病学}{确定人群中疾病或其他健康结局频率及决定因素的科学。例如，考察内城区手枪暴力的流行病学。}
\glossentry{Equipoise}{临床均衡}{对随机试验哪一组更好存在真实不确定或争议，使研究者认为让随机化决定照护不会伤害参与者。例如，专家与临床人员对引导想象能否有效治疗哮喘存在均衡。}
\glossentry{Equivalence study}{等效性研究}{旨在证明两种或多种治疗结局相似的研究，通常一种新治疗、另一种已知有效。例如，用等效性设计比较新抗生素A与旧药B治疗尿路感染。}
\glossentry{Ethnography}{民族志}{记录群体、团队、组织或社区的社会互动与行为，对文化作丰富整体分析的方法。例如，跟踪外科住院医师，探索其如何学会理解不同类型的错误。
\sourcepage{692}}
\glossentry{Exclusion criteria}{排除标准}{使潜在参与者不能入组的属性清单。例如，过去2年接受过抗抑郁药、当前使用$\alpha$或$\beta$受体阻滞剂、不能按六年级水平阅读英语。另见\emph{纳入标准}。}
\glossentry{Exempt from IRB review}{免于伦理审查}{不需伦理委员会审查的研究，如调查访谈，以及既有去标识化记录、标本的二次分析。例如，调查医学生自觉对临终讨论准备如何，被认定豁免。}
\glossentry{Expedited IRB review}{伦理快速审查}{某些对参与者风险极小的研究，由一位委员而非全体委员会快速审查。例如，采鼻拭子判定病毒感染的研究获快速审查。}
\glossentry{Experiment}{实验}{见\emph{临床试验}。}
\glossentry{Exploratory study}{探索性研究}{定性研究中，对几乎未知主题，如异常临床发现或干预意外失败的研究，方案常在过程中大幅修订。例如，考察一个新的工作场所体重管理项目为何招募不足预期人数10\%。另见\emph{比较性研究、描述性研究}。}
\glossentry{Exposure}{暴露}{表示参与者具有被认为与结局相关的某种属性。例如，胃肠出血因素研究把阿司匹林暴露定义为过去6个月平均每周至少一片、任何剂量的药片。另见\emph{保护因素、危险因素}。}
\glossentry{External validity}{外部效度}{研究结论适用于研究之外的人与事件的程度。例如，稿件认为50岁以下者出血性卒中死亡较高，但样本来自专治颅内动脉瘤的四级医疗中心，评审因而质疑外部效度。与\emph{可推广性}同义。
\sourcepage{693}}
\glossentry{Fabrication}{捏造}{编造数据或结果，并把它们记录或报告出来。例如，若团队中有多人能够接触、监督原始资料和数据分析，捏造就较不容易发生。}
\glossentry{Face validity}{表面效度}{依据是否看起来合理，判断测量似乎在多大程度上反映目标现象；一般并非很可靠的效度评价方式。例如，青少年受欢迎程度工具包含研究者认为可区分受欢迎与不受欢迎学生的条目，因而被认为有表面效度。另见\emph{构念效度、内容效度、效标关联效度}。}
\glossentry{Factorial trial}{析因试验}{研究两种或多种治疗（如A、B）的临床试验，有时涉及两个不相关的结局。参与者随机接受活性A加安慰剂B、活性B加安慰剂A、两种活性治疗，或两种安慰剂。例如，研究长期服用$\beta$胡萝卜素和阿司匹林是否影响胃肠癌风险的析因试验。}
\glossentry{False-negative result}{假阴性结果}{（1）实际有目标疾病的患者，检查结果却为阴性。例如，乳腺癌已经活检证实，但乳腺影像被判正常。（2）人群存在的效应在样本中未被检出，即研究结果不显著。例如，后续证明吸烟增加卒中风险，但早期病例对照未见显著作用（$P=0.23$），属于假阴性。}
\glossentry{False-positive result}{假阳性结果}{（1）实际没有目标疾病的患者，检查却为阳性。例如，无乳腺癌且6年随访中也未发生者，乳腺影像被判癌。（2）人群不存在的效应在样本被检出，即研究结果显著。例如，后续显示吸烟不增加帕金森病风险，但早期病例对照却显著提示增加（$P=0.03$）。
\sourcepage{694}}
\glossentry{Falsification}{篡改}{操弄研究资料、材料、设备或程序，或更改、遗漏数据结果，使研究记录不能真实反映实际发现。例如，Andrew Wakefield 发表麻疹、腮腺炎、风疹（MMR）疫苗与儿童孤独症风险增加相关的论文，因篡改数据被撤稿。（Rao TS, Andrade C. The MMR vaccine and autism: Sensation, refutation, retraction, and fraud. Indian J Psychiatry. 2011;53(2):95-96.）}
\glossentry{Falsification tests}{证伪检验}{通过规定一些若被否定就会令人担心主结果有偏倚的假设，考察发现可靠程度的检验。例如，若研究发现非甾体抗炎药（NSAID）增加哮喘风险，可能因前列腺素作用，则可检验 NSAID 是否也增加被认为与前列腺素无关的疾病（如尿路感染），或与哮喘的关联是否呈预期剂量反应。}
\glossentry{Field}{字段}{记录在数据库某一列的属性。例如，血压研究很可能包含体重字段。}
\glossentry{Financial conflicts of interest}{经济利益冲突}{专利、股票、期权和付款等经济安排，可能导致设计实施偏倚、过度解释阳性，或不发表阴性结果。例如，某药物临床试验 PI 持有生产试验药公司的期权，便有经济利益冲突。}
\glossentry{Focus group}{焦点小组}{由少数参与者组成的集体访谈，按他们能否对研究主题提供意见来选择。例如，在医务人员和患者中开展焦点小组，探讨对医院计划禁售含糖饮料的感受。}
\glossentry{Foreign key}{外键}{数据库某表中的字段，同时是另一表的主键，使两表能够连接。例如，Exams 的 ParticipantID 是连回 Participants 的外键；Participants 中的 ProviderID 也可作为外键连至 Provider。}
\glossentry{Framework analysis}{框架分析}{把个案资料组织成表格，行为研究概念、列为个案的演绎分析方式。例如，用框架分析比较不同诊所如何应对新政策举措。
\sourcepage{695}}
\glossentry{Frequentist analysis}{频率学派分析}{常见设计与分析方法，计算研究结果检验统计量的$P$值，与预设显著性水平$\alpha$比较。依据是无效应零假设下结果出现的频率，不使用结果可能性的先验信息，如其他类似问题研究。例如，分析认为新手术使早期胰腺癌患者5年死亡率降低10\%（$P=0.003$）。另见\emph{贝叶斯分析}。}
\glossentry{Generalizability}{可推广性}{见\emph{外部效度}。}
\glossentry{Ghost authors}{幽灵作者}{对研究论文、摘要等作出重大贡献，却未列名的人，常用以隐瞒资助方在设计、分析或写作中的作用。例如，试验药开发公司一位雇员写了关键临床试验结果论文，却不在作者名单中。}
\glossentry{Gold standard}{金标准}{明确判定患者有无疾病或结局的方法。例如，髋骨折诊断的金标准，要求有专科资格的放射医生审阅该髋部全部影像后确认。}
\glossentry{Good clinical practice}{药物临床试验质量管理规范}{为需要 FDA 或其他监管机构批准的药物试验建立的高质量研究指南，定义为“设计、实施、记录和报告有人类参与的试验的国际伦理和科学质量标准”。例如，研究者通过编写完整操作手册并经常检查工作，保证人员遵守要求。}
\glossentry{Grounded theory}{扎根理论}{定性研究中用数据发展新理论的归纳方法。编码产生主题，再分析、发展成新的理论命题。例如，根据 ICU 观察，用扎根理论形成对院内死亡经历的理论理解。}
\sourcepage{696}
\glossentry{Guest author}{挂名作者}{对研究论文、摘要等只有琐小贡献，却被列为作者的人，也称荣誉署名。例如，研究者邀请领域内知名专家挂名，因为认为这样编辑更可能把稿件送外审。}
\glossentry{Hazard rate}{风险率}{衡量结局在人群中瞬时发生速率的流行病学术语。实际应用几乎总以结局发生率估计。例如，50--59岁女性发生冠状动脉病的风险率估为每年0.008。}
\glossentry{Hazard ratio}{风险率比}{暴露者风险率除以未暴露者风险率，几乎总从比例风险率（Cox）模型估计。例如，50--59岁男性相对同龄女性发生冠脉病的风险率比为2.0。}
\glossentry{Health Insurance Portability and Accountability Act (HIPAA)}{健康保险流通与责任法案}{美国保护日常医疗、收费或行政过程中收集的可识别个人信息的法律（见\emph{受保护健康信息}）。按其规定，个人须签署授权，才能在研究中使用这类资料。例如，用 EHR 的可识别资料研究阿尔茨海默病遗传前，需要患者 HIPAA 授权。}
\glossentry{Health literacy}{健康素养}{个人获取、处理并理解基本健康信息的能力。例如，顾问建议把问卷的专业词 axilla 换为日常词 armpit（腋下），把 emesis 换为 vomiting（呕吐），使一般健康素养者也能理解。}
\glossentry{Healthy user effect}{健康使用者效应}{观察研究中的混杂：进行某活动或使用某产品者，可能原本比未使用者健康，因此看起来结局更好。例如，健康者更可能使用身体活动追踪设备；即使设备本身无效，评价仍可能显示使用者心血管死亡率较低。}
\glossentry{Heterogeneity}{异质性}{预测与结局关联在不同研究或参与者亚组不一致的情况。例如，绝经后雌激素对情绪、认知的研究差异很大，有的有益，有的有害，有的无效。另见\emph{同质性}。
\sourcepage{697}}
\glossentry{Highest posterior density interval}{最高后验密度区间}{概括治疗效应后验分布中给定比例的最窄区间。技术上，是后验曲线横轴上覆盖该比例的最短距离，因此确定概率密度最高的区间。例如，减重药95\%最高后验密度区间为$-1.5$至$-4.6$ kg，效应在两数之间概率0.95，其他95\%可信区间不会更窄。}
\glossentry{Homogeneity}{同质性}{预测与结局关联在不同研究中一致的情况。例如，规模适当的吸烟与肺癌研究均发现吸烟者风险明显增加，具有同质性。另见\emph{异质性}。}
\glossentry{Hospital-based controls}{医院来源对照}{病例对照研究中，从病例来源的同一所或多所医院选择的对照。例如，加工肉与上消化道癌研究，以同院接受治疗的非恶性胃肠病患者作对照。}
\glossentry{Hypothesis}{假设}{关于预期研究发现的信念陈述总称。例如，研究假设长期使用抗癫痫药与口腔癌风险增加有关。另见\emph{零假设、研究假设}。}
\glossentry{Immortal time bias}{不死时间偏倚}{生存或其他事件时间比较未准确反映风险人时而产生的偏倚。名称来自某些设计，通常因暴露与未暴露、治疗与未治疗的定义，使某些人一段时间不可能死亡，却仍被计为生存。例如，一项研究比较老年患者的出院后死亡率，分组依据是他们是否在出院后14天内配取了出院处方药。出院至配药（如第10天）之间的时间是“不死”的，因为能在第10天配药者前9天不可能死亡。不处理不死时间偏倚，会使配药看起来比实际更有益。}
\sourcepage{698}
\glossentry{Incidence}{累积发病率}{随访期内发生结局的参与者比例，有时称发病比例或累积发病率。例如，研究发现素食孕妇早产的累积发病率低于吃肉孕妇。}
\glossentry{Incidence rate}{发病密度}{原先没有该疾病的一组人中，疾病或结局发生的速率。通常以新病例数除以风险人时。例如，中年男性心梗发病密度每1000人年35例，约为中年女性17例的两倍。}
\glossentry{Incidence-density sampling}{发病密度抽样}{巢式病例对照中，重要暴露随时间改变，因而需使病例对照在相似时间测量时，选择对照的技术。例如，抗组胺药使用有季节性，可能经跌倒增加短期髋骨折；研究采用发病密度抽样，使对照在病例骨折同一个月测量用药。}
\glossentry{Inclusion criteria}{纳入标准}{研究要求潜在参与者具备的属性清单。例如，18--65岁、居住旧金山、没有抑郁病史。另见\emph{排除标准}。}
\glossentry{Incorporation bias}{纳入偏倚}{评价待评检验时，所用金标准包含待评检验结果，使检验显得比实际更有用的偏倚。例如，评价脂肪酶诊断胰腺炎时，使用的共识金标准中就有脂肪酶异常这一条。}
\glossentry{Independent}{独立}{（1）两变量不相互影响。例如，研究未见膳食坚果摄入影响血糖或血糖影响坚果摄入，因此认为两者独立。（2）某变量对另一变量的作用不依赖第三变量。例如，因为母亲教育与母乳喂养相关，研究者调整教育，估计母乳喂养对2岁语言技能的独立效应。
\sourcepage{699}}
\glossentry{Independent variable}{自变量}{见\emph{预测变量}。}
\glossentry{Index test}{待评价检验}{诊断或预后检验研究中正在评价的检验。诊断时与金标准比较，预后时与参与者实际发生情况比较。例如，腹部超声与病理金标准比较，评价阑尾炎诊断。}
\glossentry{Induction}{归纳}{通过审阅实证资料，自下而上形成理论或假设的分析方法。例如，观察住院患者，形成关于死亡经历的新理论。另见\emph{演绎}。}
\glossentry{Inference}{推断}{根据样本观察形成结论的过程。例如，调整吸烟、染料暴露等其他已知膀胱癌原因后，报告饮用井水的膀胱癌病例人数为对照的两倍（$P=0.02$），研究者据此推断饮井水使风险加倍。}
\glossentry{Informed consent}{知情同意}{告知潜在参与者研究主要内容、参加后需做什么、风险收益和替代选择，再取得其参加同意的过程。例如，伦理委员会因同意书含糊、未披露重要治疗风险而要求修订。}
\glossentry{Institutional review board (IRB)}{伦理审查委员会}{审查研究提案，确保伦理可接受、保护参与者福祉与权利的委员会，也称独立伦理委员会、伦理评审委员会、人体研究委员会或研究伦理委员会。它有权否决或要求修改。例如，大学 IRB 因参与者隐私保护不足而拒批研究。}
\glossentry{Instrument bias}{仪器偏倚}{仪器持续高估或低估真实值造成的测量偏倚。例如，漏气使呼出气体体积持续低估。
\sourcepage{700}}
\glossentry{Instrument variability}{仪器变异}{测量仪器变化导致重复测量差异。例如，股骨近端骨密度可能随髋部放置位置不同而改变，产生仪器变异。}
\glossentry{Instrumental variable}{工具变量}{与预测变量相关，但除此之外不与结局相关的变量，因此可间接估计预测对结局的作用。例如，新流感疫苗使用存在明显地区差异，研究者便以居住地区为工具变量，研究对老年人死亡的影响。}
\glossentry{Intended sample}{预定样本}{方案中规定、研究者原打算纳入的一群人。例如，2021年1月1日至6月30日，在 Longview 医院星期一或星期四首次接受治疗、且距初诊不超过6周的乳腺癌女性；这两天是研究者或人员有空的日子。另见\emph{可接触人群、样本}。}
\glossentry{Intention-to-treat analysis}{意向治疗分析}{随机试验中按原分配组比较参与者，即使实际接受治疗不同。原书将其称为最严谨的分析形式。例如，比较随机分配接受6个月心理治疗，与随机分配获得减压宣传册的对照，对焦虑症状的改善。另见\emph{按实际治疗分析、符合方案分析}。}
\glossentry{Interaction}{交互作用}{长期与效应修饰同义使用，现在则还暗示效应修饰因素是结局的原因。例如，报告吸烟、石棉与肺癌存在交互：石棉在吸烟者中引起的风险增加大于不吸烟者。}
\glossentry{Interaction term}{交互项}{线性或 logistic 等多变量模型中，用于估计两个或更多预测变量组合所产生、超出各自作用的“额外”结局效应的项。对于0／1二分类预测最容易理解，也取决于基础风险模型。例如，在乘法风险模型下，若男性使膀胱癌风险增至2倍，吸烟使风险增至5倍，则吸烟男性的风险应为不吸烟女性的10倍（$2\times5$）。在模型中加入男性（取值1）乘吸烟（取值1）等于1的交互项，可判断男性与吸烟的组合作用实际导致高于还是低于10倍的风险。
\sourcepage{701}}
\glossentry{Intercoder reliability}{编码者间信度}{两个或多个人对资料使用相同标签或编码的程度。例如，“提到一位有帮助的工作人员”的编码者间信度，高于“描述个人困惑”。}
\glossentry{Interim monitoring}{期中监测}{试验资料收集时，间歇监测是否应为保护安全而早停或改方案。例如，心梗存活者抑制室性早搏的三种抗心律失常药对安慰剂试验，发现干预组死亡远高于对照而停止。（Echt DS, Liebson PR, Mitchell LB, et al. Mortality and morbidity in patients receiving encainide, flecainide, or placebo—the cardiac arrhythmia suppression Trial. N Engl J Med. 1991;324:781-788；The Cardiac Arrhythmia Suppression Trial II Investigators. Effect of the antiarrhythmic agent moricizine on survival after myocardial infarction. N Engl J Med. 1992;327:227-233.）}
\glossentry{Intermediate marker}{中间标志物}{与临床结局相关，但尚未证明改变它会改变该结局的测量。例如，血清 S100B 在抑郁女性高于对照，因此被发现为抑郁的中间标志物。另见\emph{中介变量、替代标志物}。}
\glossentry{Internal consistency}{内部一致性}{多条目测量中，不同条目回答彼此相关的程度。例如，若所有拟测同一总体构念的问题都得到相似分数，问卷就有高内部一致性。}
\glossentry{Internal validation}{内部验证}{在与模型开发样本相同的可接触人群中，验证模型，最常用于临床预测。例如，新生儿脓毒症模型用巢式病例对照一半病例开发，再在另一半内部验证。
\sourcepage{702}}
\glossentry{Internal validity}{内部效度}{研究结论反映研究中实际发生情况的程度。例如，开始与结束时各用自动血压计测三次，以增强“冥想不比读诗更能降低升高的收缩压”这一结论的内部效度。另见\emph{外部效度}。}
\glossentry{Interpretation}{解释}{定性研究对资料的实证分析，强调在丰富情境中整体理解社会现象。例如，分析者用判断识别访谈要点，形成患者经历的叙述。}
\glossentry{Interrupted time series design}{中断时间序列设计}{一种断点回归及机会性设计，以时间作运行变量，使治疗或暴露概率在某时点突然改变。例如，研究北加州 Kaiser Permanente 医院胆红素检测仪器重新校准的作用，发现校准后新生儿黄疸光疗使用骤降60\%。（Kuzniewicz MW, Greene DN, Walsh EM, McCulloch CE, Newman TB. Association between laboratory calibration of a serum bilirubin assay, neonatal bilirubin levels, and phototherapy use. JAMA Pediatr. 2016;170(6):557-561.）}
\glossentry{Intervention}{干预}{随机试验中参与者接受的活性治疗。例如，焦虑心理治疗试验中，干预为持照心理师每周1小时、持续6个月、强调认知行为方法的治疗。另见\emph{对照}第二义。}
\glossentry{Justice}{公正原则}{研究伦理的基本原则，要求研究收益和负担公平分配。例如，新化疗试验只在有保险支付研究检测费者中开展，使无保险者无法获得潜在先进治疗，可能违反公正原则。}
\glossentry{Kappa}{kappa}{衡量两个或多个观察者对现象发生与否的认同，超出从边缘合计所预期程度的统计量，范围$-1$（完全不一致）至1（完全一致）。例如，两位病理医生对肝活检标本是否肝硬化，kappa 为0.55。
\sourcepage{703}}
\glossentry{K-fold cross validation}{K折交叉验证}{验证可能因过拟合而高估的模型性能的方法。例如，10折先把样本分十组，分别建十个模型，每次留一组不用，再依据十个留出组的预测和实际值估计性能。}
\glossentry{Level of statistical significance}{统计学显著性水平}{见\emph{α}。}
\glossentry{Likelihood ratio}{似然比}{定量描述某诊断结果如何影响患者有目标疾病或结局可能性的术语。定义为有病者出现该结果的概率除以无病者出现该结果的概率，可记 WOWO（with over without）。原文称“positive”似然比增加患病可能性，“negative”似然比降低它。\footnote{译注：似然比本身不是有正负号的量；此处应为大于1时增加患病可能性，介于0与1时减少，等于1时不改变。}分类或连续检验各个可能结果都有相应似然比。例如，铁蛋白$\le15\ \mu\mathrm{g/L}$诊断缺铁性贫血的 LR 为52，$\ge100\ \mu\mathrm{g/L}$则0.08。（Guyatt G, Oxman AD, Ali M, Willan A, Mcllroy W, Patterson C. Laboratory diagnosis of iron deficiency anemia. J Gen Intern Med. 1992;7(2):145-153.）}
\glossentry{Likert scale}{李克特量表}{对一个问题提供一系列间距近似相同的选择，通常五级。例如，“你有多大可能再来这个急诊就医？”的答案为：非常可能、较可能、说不上可能或不可能、较不可能、非常不可能。}
\glossentry{Limited data set}{有限数据集}{HIPAA 下去标识并删除指定标识、但保留日期（如卫生系统服务日期）、年龄和邮编的数据集，可依拟使用者签署的数据使用协议，为特定批准用途提供。例如，为研究 COVID-19 对医疗利用的影响，研究者可取得保险公司有限数据集。}
\glossentry{Logistic regression model}{logistic回归模型}{估计一个或多个预测变量对二分类结局的作用（表达为比值比），并调整其他预测与混杂变量的统计技术。例如，调整年龄、血压、糖尿病后，男性卒中可能性约女性两倍，黑人约白人三倍。
\sourcepage{704}}
\glossentry{Longitudinal}{纵向}{描述随时间进行的研究或测量，在这一意义上与横断面相反。例如，开始皮质激素治疗者白内障年发病的队列就是纵向研究。}
\glossentry{Marginals}{边缘合计}{二分类或分类资料表的行列总计的总称。例如，评价两位病理医生皮肤癌诊断一致性的$2\times2$表显示，第一位判25\%标本为癌，第二位18\%。}
\glossentry{Masking}{设盲}{见\emph{盲法}。}
\glossentry{Matching}{匹配}{选择一组参与者，使其某些属性与另一组相似的过程。例如，布鲁菌病病例对照按年龄（相差不超3岁）、性别及居住县匹配。又如，在研究车祸中使用安全带对重伤或死亡风险作用的队列研究中，比较同车使用与未使用安全带的乘员，从而匹配碰撞类型、时刻、速度等。另见\emph{过度匹配}。}
\glossentry{Maturation effect}{成熟效应}{人随时间学习、进步的倾向，可能使干预后改善，即使干预本身无效。例如，因成熟效应，有无在线训练的医学生心电图判读准确度都会提高。}
\glossentry{Mean}{均值}{样本或人群中连续变量的平均值，计算为全部值的总和除以人数。例如，中年女性样本血清胆固醇均值187 mg/dL。另见\emph{中位数、标准差}。}
\glossentry{Measurement error}{测量误差}{测量精密度或准确度，或二者，未达到完美的情况。多数变量至少有一些误差，可能仅生存状态例外。例如，新生儿听力研究每周用2 kg不锈钢砝码校准体重秤，以减少测量误差。}
\glossentry{Median}{中位数}{把样本或人群分成大致等大的两半的值，即第50百分位。连续变量有少数极高低值、会过度影响均值时常用。例如，548名医生年收入中位数275,000美元。另见\emph{均值、标准差}。
\sourcepage{705}}
\glossentry{Mediation}{中介作用}{治疗或暴露通过改变另一项结局原因来产生作用的过程，第二个中间原因称中介。例如，考察收入补贴使出生体重改善，是否通过母亲营养改善实现。}
\glossentry{Mediator}{中介变量}{由目标预测因素引起、又引起结局的变量，解释预测因素至少部分如何造成结局。例如，估计肥胖对卒中总因果效应时不调整糖尿病，因为相信肥胖增加糖尿病，后者又增加卒中。}
\glossentry{Memo}{备忘录}{定性研究中识别新兴趣主题、指导持续资料收集的短文。例如，研究者写备忘录讨论为何修改方案，开始问所有人曾否看见亲人在医院去世。}
\glossentry{Memorandum of Understanding (MOU)}{谅解备忘录}{双方制定签署、规定工作范围、预期产物任务和总体时间安排的文件。例如，学术机构与合作社区组织签 MOU，明确各阶段目标和付款，为合作提供结构。}
\glossentry{Mendelian randomization}{孟德尔随机化}{利用那些决定或影响对某种危险因素或治疗的易感性的基因的随机遗传，加强因果推断的技术。例如，在具有谷胱甘肽 S-转移酶T1基因型的母亲中，母亲使用对乙酰氨基酚与儿童哮喘的关联明显更强；该酶参与对乙酰氨基酚代谢物的解毒，因此这一发现增强了两者有因果关系的可能性。（Shaheen SO, Newson RB, Ring SM, Rose-Zerilli MJ, Holloway JW, Henderson AJ. Prenatal and infant acetaminophen exposure, antioxidant gene polymorphisms, and childhood asthma. J Allergy Clin Immunol. 2010;126(6):1141-8.e7.）}
\sourcepage{706}
\glossentry{Meta-analysis}{荟萃分析}{把预测与结局变量相似的多项研究结果，合成单一汇总结果的过程。例如，九项研究荟萃发现高出生体重与哮喘风险增加约20\%相关。}
\glossentry{Meta-regression}{荟萃回归}{分析研究所选特征怎样影响其结果的统计技术。例如，分析30项他汀随机试验，考察年龄较大是否预测较低疗效，同时调整男性比例、同期阿司匹林使用率等研究特征。}
\glossentry{Minimal clinically important difference}{最小临床重要差异}{工具得分中，按其与临床结局或指标的对应判断，最小有意义的差异或变化，常在研究前确定并用于样本量估计。例如，失眠治疗使睡眠质量得分少量增加，但低于最小临床重要差异阈值。}
\glossentry{Minimum necessary data}{最低必需数据}{HIPAA 定义为完成某项特定批准研究所需的最少数据要素。例如，伦理委员会可允许未取得患者同意使用既有 EHR，但只能提取、使用该项目最低必需资料。}
\glossentry{Misclassification}{错分}{分类或二分类变量的测量错误，把原属某个值的参与者错误计为另一值。例如，病历不全使确曾住院跌倒者被归为未跌倒。另见\emph{差异性错分偏倚、非差异性错分偏倚}。}
\glossentry{Missing data}{缺失数据}{研究开始或随访期间未收集的数据，通常指有些人收集、有些人没有，而非从来未收集的资料。例如，饮酒资料缺失比例较大（34\%），使研究者担心跌倒危险因素结果有偏倚。}
\glossentry{Multiple cohort study}{多队列研究}{纳入两个或更多不同群体（队列）后比较结局，常用于职业暴露，比较者可能暴露或未暴露于危险因素。例如，研究飞行中宇宙射线暴露是否与血液恶性肿瘤风险增加相关，纳入飞行员、空乘（暴露）及售票员、登机口人员（不暴露）四队列。另见\emph{双队列研究}。
\sourcepage{707}}
\glossentry{Multiple hypothesis testing}{多重假设检验}{一项研究检验不止一个、通常远多于一个假设，因此若不调整显著性水平，就会增加I类错误。例如，研究报告维生素C补充与认知下降显著相关（$P=0.03$），但因考察30多种营养补充剂而未考虑多重检验，受到批评。另见\emph{Bonferroni校正}。}
\glossentry{Multivariable analysis}{多变量分析}{调整一个或多个潜在混杂对预测与结局关联影响的统计技术总称。例如，研究调整年龄、性别、教育、研究开始时认知功能及吸烟后，发现每日饮用至少两份酒精饮料与认知下降风险增加相关。}
